\chapter{Glossar -- Common Terminology}
\label{chap:glossar}




\par\mbox{}\par
\begin{lernzieleEnv}
Alphabetisches Nachschlagewerk für Fachbegriffe, Abkürzungen und Konzepte. Nachschlagen bei unbekannten Begriffen.
\end{lernzieleEnv}

\section{A}

\textbf{AGI (Artificial General Intelligence)}\hspace{0pt}: KI mit menschenähnlicher allgemeiner Intelligenz -- ein Modell, das beliebige intellektuelle Aufgaben gleich gut wie ein Mensch löst. Bisher existiert AGI nicht; alle Modelle sind Narrow AI.

\textbf{AI Act}\hspace{0pt}: Die EU-Verordnung 2024/1689 zur Regulierung von Künstlicher Intelligenz. Kategorisiert KI-Systeme nach Risikoklassen (minimal, limited, high, unacceptable) und definiert Konformitätsanforderungen.

\textbf{AI Red Teaming}\hspace{0pt}: Strukturierte Sicherheitstests, bei denen versucht wird, ein KI-System zu manipulieren oder Schwachstellen auszunutzen. Kombiniert automatisierte Tools (Garak, PyRIT) mit manuellen Tests.

\textbf{Alignment}\hspace{0pt}: Der Prozess, ein KI-Modell auf menschliche Werte und Intentionen auszurichten. Methoden: RLHF, Constitutional AI, DPO. Ein nicht-aligniertes Modell kann hochintelligent, aber gefährlich sein.

\textbf{API (Application Programming Interface)}\hspace{0pt}: Schnittstelle zur programmatischen Interaktion mit einem Dienst -- z.B. die OpenAI Chat API zum Senden von Prompts und Empfangen von Antworten.

\textbf{Attention}\hspace{0pt}: Der Mechanismus in Transformer-Modellen, der jedem Wort im Input eine Gewichtung gibt und so den Kontext des Modells steuert. Das Herzstück aller modernen LLMs.

\textbf{Attention-Cost}\hspace{0pt}: Die Rechenkosten der Attention-Mechanismen, die quadratisch mit der Sequenzlänge skalieren -- der Haupttreiber für Latenz und Kosten bei langen Prompts.

\textbf{Autoregressive Generation}\hspace{0pt}: Die Eigenschaft von LLMs, Tokens nacheinander zu generieren, wobei jeder neue Token auf allen vorherigen aufbaut. Daher: Output-Token sind teurer als Input-Token.

\section{B}

\textbf{Base Model / Foundation Model}\hspace{0pt}: Das vortrainierte Basismodell vor jeglicher Feinabstimmung (Fine-Tuning) oder Alignment. Beispiel: Llama 3.1 70B ist das Base Model, die darauf feinabgestimmte Variante heisst z.B. Llama 3.1 Instruct.

\textbf{BLEU (Bilingual Evaluation Understudy)}\hspace{0pt}: NLP-Metrik, die die N-Gramm-Überlappung zwischen generiertem und Referenztext misst. Ursprünglich für Maschinenübersetzung entwickelt. Schwäche: Bestraft paraphrasierte, aber korrekte Antworten.

\textbf{BM25}\hspace{0pt}: Ein klassischer Text-Retrieval-Algorithmus für Suchmaschinen -- oft in Kombination mit Embeddings (hybride Suche) in RAG-Systemen verwendet.

\textbf{BPE (Byte-Pair Encoding)}\hspace{0pt}: Die am weitesten verbreitete Tokenisierungsmethode. Zerlegt Text in Subwörter basierend auf der Häufigkeit von Zeichenpaaren. Grundlage von GPT-Tokenizern.

\section{C}

\textbf{Chain-of-Thought (CoT)}\hspace{0pt}: Eine Prompting-Technik, bei der das Modell Schritt-für-Schritt nachdenkt, bevor es antwortet -- verbessert die Qualität bei logischen und analytischen Aufgaben signifikant.

\textbf{Chat Completions API}\hspace{0pt}: OpenAIs Interface für Interaktion mit GPT-Modellen via Message-Passing (system/user/assistant). Das Standard-Interface für LLM-Integrationen.

\textbf{Chunking}\hspace{0pt}: Der Prozess, grosse Dokumente in kleinere Textabschnitte (Chunks) zu zerlegen -- der erste Schritt bei der RAG-Indexierung.

\textbf{Circuit Breaker}\hspace{0pt}: Ein Architekturmuster, das bei wiederholten Fehlern einer externen Abhängigkeit (z.B. LLM-API) automatisch weitere Anfragen blockiert und eine Fallback-Antwort liefert.

\textbf{Claude}\hspace{0pt}: Anthropics Familie von LLMs (Opus, Sonnet, Haiku). Die Alternative zu OpenAIs GPT-Reihe. Die konkreten Modellversionen ändern sich regelmässig.

\textbf{Context Window / Kontextfenster}\hspace{0pt}: Die maximale Anzahl von Tokens, die ein Modell auf einmal verarbeiten kann -- je nach Modell zwischen 8K und 1M Tokens.

\textbf{Constitutional AI}\hspace{0pt}: Anthropics Methode für Alignment, bei dem das Modell gegen eine schriftliche Verfassung (Constitution) trainiert wird. Das Modell lernt, eigene Outputs gegen diese Prinzipien zu prüfen und zu korrigieren.

\textbf{Cosine Similarity}\hspace{0pt}: Ein Mass für die Ähnlichkeit zwischen zwei Vektoren (Embeddings). Werte nahe 1 bedeuten hohe Ähnlichkeit, Werte nahe 0 keine. Die Standard-Metrik für semantische Suche in Vector Databases.

\textbf{Cutoff Date}\hspace{0pt}: Das Datum, bis zu dem ein Modell im Training Daten gesehen hat. Ein Modell kennt nur Daten bis zu seinem Cutoff-Datum. Ereignisse danach sind ihm unbekannt.

\section{D}

\textbf{Data Poisoning}\hspace{0pt}: Ein Angriff, bei dem manipulierte Daten in den Trainingsdatensatz eines Modells eingeschleust werden, um sein Verhalten gezielt zu beeinflussen.

\textbf{Decoder}\hspace{0pt}: Der Teil eines Transformer-Modells, der Text generiert (im Gegensatz zum Encoder, der Text repräsentiert). Decoder-only-Modelle (GPT, Claude, Llama) sind der aktuelle Standard.

\textbf{Deep Learning}\hspace{0pt}: Ein Teilbereich des ML, der neuronale Netze mit vielen Schichten verwendet -- die Grundlage aller modernen LLMs.

\textbf{DPO (Direct Preference Optimization)}\hspace{0pt}: Eine Alternative zu RLHF für Alignment. Statt eines Reward-Modells wird direkt auf menschlichen Präferenzen trainiert. Einfacher, stabiler und oft genauso effektiv.

\section{E}

\textbf{Embedding}\hspace{0pt}: Ein Vektor (Liste von Floats), in den ein Text durch ein Embedding-Modell umgewandelt wird. Ähnliche Texte haben ähnliche Embeddings -- die Grundlage semantischer Suche.

\textbf{EOS (End of Sequence)}\hspace{0pt}: Das spezielle Token, das signalisiert, dass eine Generation endet -- das Modell hat seine Antwort beendet.

\textbf{EU AI Act}\hspace{0pt}: Siehe AI Act.

\textbf{Eval / Evaluation}\hspace{0pt}: Der Prozess, die Qualität eines LLM-Systems systematisch zu messen. Umfasst Golden Datasets, LLM-as-a-Judge, Metriken und Monitoring.

\section{F}

\textbf{Few-Shot Prompting}\hspace{0pt}: Eine Prompting-Technik, bei der dem Modell mehrere Beispiel-Eingabe-Ausgabe-Paare mitgegeben werden -- der Standard für konsistente Produktivsysteme.

\textbf{FIM (Fill-in-the-Middle)}\hspace{0pt}: Eine Fähigkeit von Code-LLMs, mitten in einem Dokument zu generieren statt nur am Ende. Wichtig für IDEs wie GitHub Copilot.

\textbf{Fine-Tuning}\hspace{0pt}: Das Anpassen eines vortrainierten Modells auf eine spezifische Aufgabe oder Domäne -- durch zusätzliches Training auf eigenen Daten, ohne das gesamte Modell von null zu trainieren.

\textbf{Foundation Model}\hspace{0pt}: Siehe Base Model.

\textbf{Function Calling / Tool Use}\hspace{0pt}: Die Fähigkeit eines LLMs, strukturierte Funktionsaufrufe zu generieren, die dann von einem externen System ausgeführt werden. Grundlage jeder Agenten-Architektur.

\section{G}

\textbf{Garak}\hspace{0pt}: Open-Source-Tool für automatisierte LLM-Sicherheitstests. Prüft ein Modell systematisch auf Jailbreak-Resistenz, Halluzinationen, Toxizität und weitere Risikokategorien.

\textbf{Gemini}\hspace{0pt}: Googles Familie von LLMs. Besonders stark in Multimodalität und sehr grossen Kontextfenstern.

\textbf{Golden Dataset}\hspace{0pt}: Eine Sammlung von Input-Output-Paaren mit bekannten, korrekten Antworten -- der Standard für reproduzierbare Prompt-Evaluation.

\textbf{GPT (Generative Pre-trained Transformer)}\hspace{0pt}: Die von OpenAI entwickelte Modellfamilie. Die jeweils aktuelle Version ist das meistgenutzte Multimodal-Modell des Anbieters.

\textbf{GPU (Graphics Processing Unit)}\hspace{0pt}: Ursprünglich für Grafikberechnung entwickelt, heute das Standard-Hardware für KI-Training und -Inference. NVIDIA H100/B200 sind die führenden KI-GPUs.

\textbf{Guardrails}\hspace{0pt}: Ein Framework oder eine Schicht, die die Ausgabe eines LLMs gegen definierte Regeln validiert und bei Verstössen eingreift (blockieren, umformulieren, eskalieren).

\section{H}

\textbf{Halluzination}\hspace{0pt}: Wenn ein LLM falsche oder erfundene Informationen als Fakten ausgibt. Das Modell erfindet Inhalte mit absolutem Selbstvertrauen -- kein Bug, sondern inhärente Eigenschaft generativer Modelle.




\par\mbox{}\par
\begin{praxishinweisEnv}
Halluzinationen lassen sich nie vollständig eliminieren -- nur reduzieren. Bewährte Strategien: RAG mit Quellenangabe, Golden Dataset für Regressionstests, LLM-as-a-Judge für Output-Validierung. Bei sicherheitskritischen Anwendungen: Human-in-the-Loop einplanen.
\end{praxishinweisEnv}

\textbf{HITL (Human-in-the-Loop)}\hspace{0pt}: Ein Workflow, bei dem ein Mensch bei kritischen Entscheidungen eingreift oder die finale Freigabe gibt. Pflicht bei Hochrisiko-KI-Systemen.

\textbf{Hybrid Search}\hspace{0pt}: Kombination aus Keyword-Suche (BM25) und semantischer Suche (Embeddings) in Vector Databases -- löst die Schwächen der reinen Ansätze auf.

\section{I}

\textbf{Inference}\hspace{0pt}: Der Prozess, bei dem ein trainiertes Modell eine Vorhersage trifft oder Text generiert -- im Gegensatz zum Training. Inference ist der Betriebsmodus in der Produktion.

\textbf{Instruction Tuning}\hspace{0pt}: Das Fine-Tuning eines Base Models auf Anweisungsbefolgung (Instruction-Following). Der Schritt, der aus einem Base Model ein Chat-Modell macht.

\textbf{IOU (Input/Output Unit)}\hspace{0pt}: Kosteneinheit bei LLM-APIs. Input-Tokens sind günstiger (ca. Faktor 3--5) als Output-Tokens.

\section{J}

\textbf{Jailbreak}\hspace{0pt}: Ein Angriff, der die Safety-Restriktionen eines LLMs umgeht. Beispiele: "Do Anything Now" (DAN), Rollenspiel-Jailbreaks, Payload-Splitting.

\textbf{JSON Mode}\hspace{0pt}: Ein API-Modus, der das LLM zwingt, gültiges JSON auszugeben. Ermöglicht strukturierte Outputs ohne Parser-Fehler.

\section{K}

\textbf{Knowledge Cutoff}\hspace{0pt}: Siehe Cutoff Date.

\textbf{KVCache (Key-Value Cache)}\hspace{0pt}: Ein Cache-Mechanismus in Transformer-Modellen, der die berechneten Key-Value-Paare für bereits generierte Tokens speichert. Spart Rechenzeit bei autoregressiver Generation.

\section{L}

\textbf{LangChain / LangGraph}\hspace{0pt}: Framework zur Entwicklung von LLM-Anwendungen. LangChain fokussiert auf Chains und RAG, LangGraph auf zustandsbehaftete Multi-Agent-Workflows.

\textbf{Langfuse}\hspace{0pt}: Open-Source-Plattform für Tracing, Monitoring und Evaluation von LLM-Anwendungen. Ermöglicht detaillierte Einblicke in Kosten, Latenz und Qualität.

\textbf{LLM (Large Language Model)}\hspace{0pt}: Ein neuronales Netz, das auf riesigen Textkorpora trainiert wurde, um natürliche Sprache zu verstehen und zu generieren. Wichtigste aktuelle Vertreter: GPT (OpenAI), Claude (Anthropic), Gemini (Google), Llama (Meta).

\textbf{LLM-as-a-Judge}\hspace{0pt}: Eine Evaluationsmethode, bei der ein LLM die Qualität der Antwort eines anderen LLMs bewertet. Günstig und skalierbar, aber anfällig für Bias (Position, Länge, Stil).

\textbf{Logit}\hspace{0pt}: Die rohe Log-Wahrscheinlichkeit eines Tokens vor der Anwendung von Softmax (Temperature). Kann für Confidence-Scoring verwendet werden.

\textbf{LoRA (Low-Rank Adaptation)}\hspace{0pt}: Eine Technik für parameter-effizientes Fine-Tuning: Statt aller Modell-Gewichte werden nur kleine Adapter-Matrizen trainiert -- spart bis zu 95\% der Trainingskosten.

\textbf{LSTM}\hspace{0pt}: Eine Architektur von rekurrenten neuronalen Netzen (RNN) -- Vorgänger der Transformer-Architektur.

\section{M}

\textbf{MCP (Model Context Protocol)}\hspace{0pt}: Ein offenes Protokoll von Anthropic, das LLMs mit externen Tools und Datenquellen verbindet. Alternativ zu Function Calling.

\textbf{MHA (Multi-Head Attention)}\hspace{0pt}: Die Variante der Attention in Transformatoren, bei der mehrere Heads parallel arbeiten und unterschiedliche Aspekte des Kontexts erfassen. Standard in allen modernen LLMs.

\textbf{ML (Machine Learning)}\hspace{0pt}: Maschinelles Lernen.

\textbf{MoE (Mixture of Experts)}\hspace{0pt}: Eine Architektur, bei der das Modell mehrere kleinere Experten-Netzwerke hat und nur die relevantesten pro Token aktiviert -- spart Rechenkosten bei grossen Modellen.

\textbf{Multimodalität}\hspace{0pt}: Die Fähigkeit eines Modells, mehrere Datentypen zu verarbeiten (Text, Bilder, Audio, Video). Die Flaggschiff-Modelle von OpenAI, Google und zunehmend Anthropic sind multimodal.

\section{N}

\textbf{NER (Named Entity Recognition)}\hspace{0pt}: Die Aufgabe, benannte Entitäten (Personen, Orte, Organisationen) in Text zu identifizieren. Ein klassischer NLP-Task, den LLMs gut lösen.

\textbf{NGramm}\hspace{0pt}: Eine Folge von N aufeinanderfolgenden Token in einem Text. Grundlage vieler klassischer NLP-Metriken (BLEU, ROUGE).

\textbf{NPU (Neural Processing Unit)}\hspace{0pt}: Spezialisierte Hardware für KI-Inference auf Endgeräten. Apple Neural Engine und Qualcomm Hexagon NPU sind Beispiele.

\section{O}

\textbf{OpenAI}\hspace{0pt}: Das Unternehmen hinter GPT, DALL-E, Whisper und der OpenAI-Plattform. Der meistgenutzte LLM-Anbieter weltweit.

\textbf{Output-Token}\hspace{0pt}: Das vom Modell generierte Token -- fast immer teurer als Input-Token (3--5x Faktor) und damit ein kritischer Kostenfaktor.

\textbf{OWASP LLM Top 10}\hspace{0pt}: Ein Sicherheitskatalog, der die zehn wichtigsten Risikokategorien für LLM-Anwendungen auflistet. Standard-Referenz für LLM-Sicherheit.

\section{P}

\textbf{Parameter}\hspace{0pt}: Die internen Gewichte eines neuronalen Netzes, die durch Training gelernt werden. Je mehr Parameter, desto komplexere Muster kann ein Modell lernen -- aber desto teurer ist die Inference.

\textbf{PII (Personally Identifiable Information)}\hspace{0pt}: Personenbezogene Daten wie Namen, E-Mails, Telefonnummern, Adressen. Dürfen nicht ungeschützt an LLM-APIs gesendet werden (DSGVO).

\textbf{Pinecone}\hspace{0pt}: Eine Managed Vector Database as a Service. Einfachster Weg, Embeddings in Produktion zu skalieren.

\textbf{Prompt}\hspace{0pt}: Die Eingabe eines Users an ein LLM. Im weiteren Sinne umfasst der Prompt auch System-Prompts, Few-Shot-Beispiele und Kontext-Chunks.

\textbf{Prompt Caching}\hspace{0pt}: Ein Feature von OpenAI und Anthropic, bei dem lange System-Prompts zwischengespeichert werden -- nach dem ersten Aufruf zahlt man nur die Differenz (oft 90\% Rabatt).

\textbf{Prompt Injection}\hspace{0pt}: Ein Angriff, bei dem bösartiger Input in den Prompt eingeschleust wird, um das Verhalten des LLMs zu manipulieren. Die kritischste Sicherheitslücke bei LLM-Anwendungen.

\textbf{Promptfoo}\hspace{0pt}: CLI-Tool für systematische Prompt-Evaluation. Ermöglicht Golden-Dataset-Tests, A/B-Comparisons und Regression-Checks.

\textbf{Pydantic}\hspace{0pt}: Python-Bibliothek für Datenvalidierung via Type Hints. Grundlage für strukturierte Outputs in LLM-Anwendungen.

\section{Q}

\textbf{Qwen}\hspace{0pt}: Chinesische LLM-Familie von Alibaba (ehemals Tongyi Qianwen). Besonders stark in Multilingualität und Code-Generation.

\section{R}

\textbf{RAG (Retrieval-Augmented Generation)}\hspace{0pt}: Eine Architektur, bei der ein LLM vor der Antwort-Generierung relevante Dokumente aus einer externen Wissensdatenbank zieht -- das Standardmuster für Firmenwissen-Integration.




\par\mbox{}\par
\begin{praxishinweisEnv}
RAG ist meist der bessere erste Schritt als Fine-Tuning. Erlaubt Wissensaktualisierung ohne Retraining, reduziert Halluzinationen durch Quellenbindung. Faustregel: RAG zuerst, Fine-Tuning nur für Stil- oder Format-Anpassungen.
\end{praxishinweisEnv}

\textbf{Rate Limiting}\hspace{0pt}: Die Begrenzung der Anzahl von API-Requests pro Zeitfenster. Schützt vor Kosten-Explosion und Abuse.

\textbf{ReAct (Reasoning + Acting)}\hspace{0pt}: Ein Agent-Pattern, bei dem das Modell zwischen Reasoning (Thought) und Action (Tool Call) abwechselnd durchläuft bis zur Lösung der Aufgabe.

\textbf{Red Teaming}\hspace{0pt}: Der professionelle Test von KI-Systemen auf Sicherheitsprobleme -- versucht, das Modell zu manipulieren oder Schwachstellen zu finden.

\textbf{RLHF (Reinforcement Learning from Human Feedback)}\hspace{0pt}: Eine Alignment-Methode, bei der ein Reward-Modell auf menschlichen Präferenzen trainiert wird, um das LLM zu steuern.

\textbf{ROUGE (Recall-Oriented Understudy for Gisting Evaluation)}\hspace{0pt}: NLP-Metrik, die den Recall der N-Gramm-Überlappung misst. Verbreitet für Zusammenfassungen.

\section{S}

\textbf{Sampling}\hspace{0pt}: Der Prozess der Token-Auswahl während der Generierung. Parameter wie Temperature, Top-K und Top-P steuern die Zufälligkeit.

\textbf{SFT (Supervised Fine-Tuning)}\hspace{0pt}: Das Fine-Tuning eines Modells auf gelabelten Beispieldaten -- der Standardweg für Instruct-Modelle.

\textbf{Speculative Decoding}\hspace{0pt}: Eine Technik, bei der ein kleines Modell eine Antwort vorhersagt und ein grosses Modell diese validiert. Spart Rechenzeit bei langen Generationen.

\textbf{SSE (Server-Sent Events)}\hspace{0pt}: Ein Standard für Streaming-Responses. Das Backend sendet fortlaufend Daten-Events an den Client -- die Grundlage von Streaming in LLM-APIs.

\textbf{Structured Output / JSON Mode}\hspace{0pt}: Die Fähigkeit eines LLMs, Ausgaben in einem vorgegebenen Format (JSON, Pydantic-Modell) zu liefern. Standard für API-Integration.

\textbf{System-Prompt}\hspace{0pt}: Ein vorangestellter Instruction-Text, der das Verhalten des LLMs für die gesamte Sitzung festlegt. Das zentrale Steuerungswerkzeug des AI Engineers.

\section{T}

\textbf{Temperature}\hspace{0pt}: Ein Parameter, der die Kreativität der LLM-Generierung steuert. 0 = weitgehend deterministisch (reduziert Varianz drastisch, garantiert aber keine Bit-Identität -- serverseitiges Batching kann die Floating-Point-Arithmetik minimal verschieben), 1 = kreativ/variiert, >1 = extrem volatil.

\textbf{Token}\hspace{0pt}: Die kleinste Einheit, die ein LLM verarbeitet. Ein Token kann ein Wort, ein Teilwort oder ein einzelnes Zeichen sein.

\textbf{Tokenizer}\hspace{0pt}: Das Modul, das Text in Tokens zerlegt (Encoding) und Tokens zurück in Text (Decoding). Jedes Modell hat seinen eigenen Tokenizer.

\textbf{TTFT (Time to First Token)}\hspace{0pt}: Die Zeit bis zum ersten generierten Token. Die wichtigste Latenz-Metrik für LLM-Streaming.

\textbf{Tiktoken}\hspace{0pt}: OpenAIs Tokenisierungs-Bibliothek für Python. Das Standardwerkzeug, um Tokens zu zählen und damit Kosten und Kontext-Limits zu kontrollieren.

\textbf{Top-K Sampling}\hspace{0pt}: Eine Sampling-Methode, bei der das Modell nur aus den K wahrscheinlichsten nächsten Tokens wählt.

\textbf{Top-P (Nucleus) Sampling}\hspace{0pt}: Eine Sampling-Methode, die aus dem Cumulative-Probability-P der wahrscheinlichsten Tokens wählt. Standard in vielen Modellen.

\textbf{Tracing}\hspace{0pt}: Die Aufzeichnung eines Requests über verteilte Systemkomponenten hinweg. Jeder Schritt (API-Call, DB-Query, LLM-Call) wird als Span in einer Trace aufgezeichnet.

\textbf{Transformer}\hspace{0pt}: Die grundlegende Architektur hinter allen modernen LLMs. Basiert auf der Attention-Architektur (Vaswani et al., "Attention is All You Need", 2017).

\section{U}

\textbf{Uv}\hspace{0pt}: Moderner Python-Package-Manager (Alternative zu pip). Extrem schnell, geschrieben in Rust. der empfohlene Paketmanager für neue KI-Projekte.

\section{V}

\textbf{Vector Database}\hspace{0pt}: Eine Datenbank, die speziell für die Speicherung und Abfrage von Vektoren (Embeddings) optimiert ist -- z.B. Pinecone, ChromaDB, pgvector, Qdrant, Weaviate. Kernkomponente jeder RAG-Architektur.

\textbf{Vector Search (Ähnlichkeitssuche)}\hspace{0pt}: Die Suche nach den nächsten Nachbarn (Nearest Neighbors) in einem Vektorraum. Ermöglicht semantische Suche.

\section{W}

\textbf{Weight}\hspace{0pt}: Siehe Parameter.

\textbf{Whisper}\hspace{0pt}: OpenAIs Open-Source-Spracherkennungsmodell. Standard für Speech-to-Text in LLM-Pipelines.

\textbf{White-Box Test}\hspace{0pt}: Ein Test, bei dem der interne Zustand des Modells beobachtet wird (z.B. Token-Wahrscheinlichkeiten, Attention-Maps). Im Gegensatz zu Black-Box-Tests, die nur Input/Output betrachten.

\section{Z}

\textbf{Zero-Shot Prompting}\hspace{0pt}: Eine Technik, bei der das Modell keine Beispiele bekommt, sondern nur eine direkte Anweisung -- der einfachste, aber auch variabelste Prompt-Ansatz.

\section*{Abkürzungsverzeichnis}

\begin{longtable}{@{}p{2.5cm}p{\dimexpr\linewidth-2.5cm-2\tabcolsep}@{}}
\toprule
\textbf{Abkürzung} & \textbf{Vollständig} \\
\midrule
\endhead
AGI & Artificial General Intelligence \\
API & Application Programming Interface \\
BPE & Byte-Pair Encoding \\
CoT & Chain-of-Thought \\
DPO & Direct Preference Optimization \\
DSGVO & Datenschutz-Grundverordnung (GDPR) \\
EOS & End of Sequence \\
FIM & Fill-in-the-Middle \\
GPU & Graphics Processing Unit \\
HITL & Human-in-the-Loop \\
LLM & Large Language Model \\
LoRA & Low-Rank Adaptation \\
MCP & Model Context Protocol \\
MHA & Multi-Head Attention \\
ML & Machine Learning \\
MoE & Mixture of Experts \\
NER & Named Entity Recognition \\
NPU & Neural Processing Unit \\
PII & Personally Identifiable Information \\
RAG & Retrieval-Augmented Generation \\
RLHF & Reinforcement Learning from Human Feedback \\
SFT & Supervised Fine-Tuning \\
SLA & Service Level Agreement \\
SOTA & State of the Art \\
SSE & Server-Sent Events \\
TTFT & Time to First Token \\
\bottomrule
\caption{Wichtige Abkürzungen und ihre Bedeutungen}
\end{longtable}
\endinput
